\documentclass[12pt]{article}
\pdfoutput=1

\usepackage{amsfonts}
\usepackage{amsmath}
\numberwithin{equation}{section}
\usepackage{amssymb}
\usepackage{array}
\usepackage{booktabs}
\usepackage{cite}
\usepackage{enumitem}
\usepackage[margin=2.5cm]{geometry}
\usepackage{mathrsfs}
\usepackage{microtype}
\usepackage{setspace}
\usepackage{titlesec}

\setlength{\parskip}{1.5ex}
\setlength{\parindent}{0pt}
\def\baselinestretch{1.15}
\pagestyle{plain}

\usepackage{color}
\definecolor{darkblue}{rgb}{0,0,0.6}
\definecolor{purple}{rgb}{0.4,0.2,0.7}
\usepackage[bookmarks=false,colorlinks=true,linkcolor=darkblue,
  citecolor=purple,urlcolor=darkblue]{hyperref}

\renewcommand{\i}{\mathrm{i}}
\renewcommand{\d}{\mathrm{d}}
\newcommand{\id}{\mathbf{1}}
\newcommand{\cO}{\mathcal{O}}
\newcommand{\QB}{Q_{\mathrm B}}
\newcommand{\jB}{j_{\mathrm B}}
\newcommand{\TF}{T_{F}}
\newcommand{\Tmat}{T_{\mathrm m}}
\newcommand{\Ttot}{T_{\mathrm{tot}}}
\newcommand{\XPCO}{\mathcal{X}}

\title{Worldsheet Conventions for Type IIB String Theory}
\author{Raghu Mahajan}
\date{August 31, 2026}

\begin{document}

\maketitle

\tableofcontents
\clearpage

\section{Flat-space worldsheet conventions}
\label{sec:worldsheet-conventions}

On each worldsheet coordinate chart, the Euclidean worldsheet coordinates
are denoted by $z$ and $\bar z$.
Holomorphic fields depend on $z$, and anti-holomorphic fields depend on
$\bar z$.  Target-space vector indices $M,N$ take values in
$\{0,1,\ldots,9\}$.  Repeated target-space indices are summed over this set.
We use the ten-dimensional Minkowski metric
\begin{equation}
\eta_{MN}:=\operatorname{diag}(-1,+1,\ldots,+1).
\label{eq:minkowski-metric}
\end{equation}
Indices on target-space vectors are raised and lowered with
$\eta_{MN}$ and its inverse $\eta^{MN}$.  We set
\begin{equation}
\alpha'=1.
\label{eq:alpha-prime}
\end{equation}

For two holomorphic operators $A(z)$ and $B(w)$, the notation
$A(z)B(w)\sim C(z,w)$ means that, in the expansion valid for $|z|>|w|$,
$C(z,w)$ contains precisely the terms that are singular as $z\to w$.  Write
$\operatorname{Sing}_{z\to w}[A(z)B(w)]$ for the sum of those singular terms.
Every coincident-point product in these notes is understood to be normal
ordered by point splitting.  Thus the notation $A(w)B(w)$ means
\begin{equation}
A(w)B(w)
:=
\lim_{z\to w}
\left(A(z)B(w)-\operatorname{Sing}_{z\to w}[A(z)B(w)]\right).
\label{eq:normal-ordering-definition}
\end{equation}
Normal-ordering colons are suppressed everywhere below.
When a composite field is defined by an equation whose left-hand side is
$C(z)$, every field on the right-hand side is evaluated at $z$ unless a
different argument is displayed explicitly.
If $T$ is a stress tensor, a holomorphic operator $\cO_h$ has conformal
weight $h$ when
\begin{equation}
T(z)\cO_h(w)
\sim
\frac{h\,\cO_h(w)}{(z-w)^2}
+\frac{\partial\cO_h(w)}{z-w}.
\label{eq:conformal-weight-definition}
\end{equation}
The central charge $c$ of $T$ is defined so that the coefficient of
$(z-w)^{-4}$ in $T(z)T(w)$ is $c/2$.  A field of Grassmann parity
$|A|\in\{0,1\}$ is called even when $|A|=0$ and odd when $|A|=1$; away from
operator-product singularities, two such fields obey
$AB=(-1)^{|A||B|}BA$.  Derivatives without a bar mean
$\partial:=\partial/\partial z$.

\subsection{Conventions for the free boson}
\label{sec:free-boson}

Let $X^M(z,\bar z)$ be the ten embedding coordinates.  Their two-point
operator product and its differentiated holomorphic form are
\begin{align}
X^M(z,\bar z)X^N(w,\bar w)
&\sim
-\frac{\eta^{MN}}{2}\log |z-w|^2,
\label{eq:XX-ope}
\\
\partial X^M(z)\partial X^N(w)
&\sim
-\frac{\eta^{MN}}{2(z-w)^2}.
\label{eq:dX-dX-ope}
\end{align}
Equation~\eqref{eq:dX-dX-ope} is the free-boson normalization used in
Ref.~\cite{Barman:2026qzi}.  The holomorphic stress tensor of the ten bosons
is defined by
\begin{equation}
T_X(z):=-\eta_{MN}\partial X^M\partial X^N(z).
\label{eq:TX-definition}
\end{equation}
With Eqs.~\eqref{eq:dX-dX-ope} and \eqref{eq:TX-definition}, each $X^M$ has
holomorphic conformal weight zero, each $\partial X^M$ has holomorphic
conformal weight one, and the ten-boson central charge is
\begin{equation}
c_X=10.
\label{eq:cX}
\end{equation}

For a target-space momentum $k_M\in\mathbb{R}^{1,9}$, define
\begin{align}
k\mathbin{\cdot}X
&:=k_MX^M,
\label{eq:k-dot-X}
\\
k^2
&:=\eta^{MN}k_Mk_N,
\label{eq:k-squared}
\\
V_k(z,\bar z)
&:=\exp\!\bigl(\i k\mathbin{\cdot}X(z,\bar z)\bigr).
\label{eq:momentum-vertex}
\end{align}
The holomorphic and anti-holomorphic conformal weights of $V_k$ are
\begin{equation}
(h_k,\bar h_k)=\left(\frac{k^2}{4},\frac{k^2}{4}\right).
\label{eq:momentum-vertex-weight}
\end{equation}
Two useful consequences of Eq.~\eqref{eq:XX-ope} are
\begin{align}
\partial X^M(z)V_k(w,\bar w)
&\sim
-\frac{\i k^M}{2(z-w)}V_k(w,\bar w),
\label{eq:dX-V-ope}
\\
V_k(z,\bar z)V_p(w,\bar w)
&\sim
|z-w|^{k\cdot p}
\exp\!\bigl(\i(k+p)\mathbin{\cdot}X(w,\bar w)\bigr),
\label{eq:V-V-ope}
\end{align}
where $k\cdot p:=\eta^{MN}k_Mp_N$.

\subsection{Conventions for the free fermion}
\label{sec:free-fermion}

Let $\psi^M(z)$ and $\widetilde\psi^M(\bar z)$ be the holomorphic and
anti-holomorphic components of the ten real worldsheet fermions.  Their
normalization is
\begin{equation}
\psi^M(z)\psi^N(w)
\sim
-\frac{\eta^{MN}}{2(z-w)}.
\label{eq:psi-psi-ope}
\end{equation}
The holomorphic fermion stress tensor is
\begin{equation}
T_\psi(z):=\eta_{MN}\psi^M\partial\psi^N(z).
\label{eq:Tpsi-definition}
\end{equation}
Thus $\psi^M$ has holomorphic conformal weight $1/2$, and the ten-fermion
central charge is
\begin{equation}
c_\psi=5.
\label{eq:cpsi}
\end{equation}
The free matter stress tensor and the matter supercurrent are defined by
\begin{align}
\Tmat
&:=T_X+T_\psi,
\label{eq:Tm-free-definition}
\\
\TF
&:=-\eta_{MN}\psi^M\partial X^N.
\label{eq:TF-free-definition}
\end{align}
The minus sign in Eq.~\eqref{eq:TF-free-definition} is fixed jointly by
Eqs.~\eqref{eq:dX-dX-ope} and \eqref{eq:psi-psi-ope}.  It ensures that the
picture-changing operator defined in Eq.~\eqref{eq:pco-definition} maps the
$-1$-picture factor $e^{-\phi}\psi^M$ to $\partial X^M$ with positive
coefficient.  The superconformal operator products of $\Tmat$ and $\TF$ are
derived in Sec.~\ref{sec:super-virasoro}.

We next fix the gamma matrices and the Ramond spin fields.  Spinor indices
$\alpha,\beta,\gamma$ take values in $\{1,\ldots,16\}$.  We use symmetric
chiral gamma-matrix blocks
\begin{equation}
(\Gamma^M)_{\alpha\beta}=(\Gamma^M)_{\beta\alpha},
\qquad
(\Gamma^M)^{\alpha\beta}=(\Gamma^M)^{\beta\alpha}.
\label{eq:gamma-symmetry}
\end{equation}
Their normalization is defined by
\begin{equation}
(\Gamma^M)_{\alpha\beta}(\Gamma^N)^{\beta\gamma}
+(\Gamma^N)_{\alpha\beta}(\Gamma^M)^{\beta\gamma}
=2\eta^{MN}\delta_\alpha{}^\gamma.
\label{eq:chiral-clifford}
\end{equation}
Equivalently, the $32\times32$ matrices
\begin{equation}
\boldsymbol\Gamma^M
:=
\begin{pmatrix}
0&(\Gamma^M)^{\alpha\beta}\\
(\Gamma^M)_{\alpha\beta}&0
\end{pmatrix}
\label{eq:dirac-gamma-from-blocks}
\end{equation}
satisfy
\begin{equation}
\{\boldsymbol\Gamma^M,\boldsymbol\Gamma^N\}
=2\eta^{MN}\id_{32}.
\label{eq:dirac-clifford}
\end{equation}
For $p\geq1$, the antisymmetrized product is
\begin{equation}
\boldsymbol\Gamma^{M_1\cdots M_p}
:=\boldsymbol\Gamma^{[M_1}\cdots\boldsymbol\Gamma^{M_p]},
\label{eq:antisymmetrized-gamma}
\end{equation}
where the brackets include the factor $1/p!$.  A symbol such as
$(\Gamma^{M_1\cdots M_p})_\alpha{}^\beta$ denotes the indicated chiral block
of Eq.~\eqref{eq:antisymmetrized-gamma}.  We choose the target-space
orientation and positive-chirality block by
\begin{equation}
\epsilon^{01\cdots9}=+1,
\qquad
(\Gamma^{01\cdots9})_\alpha{}^\beta
=\delta_\alpha{}^\beta.
\label{eq:gamma-orientation}
\end{equation}
Equations~\eqref{eq:gamma-symmetry} and \eqref{eq:gamma-orientation} are the
ten-dimensional conventions of Ref.~\cite{Alexandrov:2021dyl}.

Let $S_\alpha(z)$ and $\widetilde S_\alpha(\bar z)$ denote the left- and
right-moving positive-chirality Ramond spin fields.  Type IIB means that
these two fields have the same ten-dimensional chirality.  Each has
conformal weight $5/8$ in its chiral sector.  The symbols $S^\alpha$ and
$\widetilde S^\alpha$ denote the spin fields that occur in the $-3/2$ picture;
the placement of the spinor index is part of this convention and is fixed by
the operator products below, rather than by an unstated index-raising rule.

The Ramond operator products also contain the bosonized superghost $\phi$.
For use in the next three equations, $\phi$ is defined by
\begin{equation}
\partial\phi(z)\partial\phi(w)
\sim-\frac{1}{(z-w)^2},
\qquad
e^{q_1\phi}(z)e^{q_2\phi}(w)
\sim(z-w)^{-q_1q_2}e^{(q_1+q_2)\phi}(w).
\label{eq:phi-provisional-definition}
\end{equation}
The complete $\beta\gamma$ bosonization is given in
Sec.~\ref{sec:beta-gamma}.  We normalize the spin fields by the singular
terms
\begin{align}
e^{-\phi}\psi^M(z)\;e^{-\phi/2}S_\alpha(w)
&\sim
\frac{\i}{2(z-w)}
(\Gamma^M)_{\alpha\beta}
e^{-3\phi/2}S^\beta(w),
\label{eq:psi-spin-ope}
\\
e^{-3\phi/2}S^\alpha(z)\;e^{-\phi/2}S_\beta(w)
&\sim
\frac{\delta^\alpha{}_\beta}{(z-w)^2}e^{-2\phi}(w),
\label{eq:spin-conjugate-ope}
\\
e^{-\phi/2}S_\alpha(z)\;e^{-\phi/2}S_\beta(w)
&\sim
\frac{\i}{z-w}(\Gamma^M)_{\alpha\beta}
e^{-\phi}\psi_M(w).
\label{eq:same-chirality-spin-ope}
\end{align}
The branch choices, phases, and factors in
Eqs.~\eqref{eq:psi-spin-ope}--\eqref{eq:same-chirality-spin-ope} are those of
Ref.~\cite{Alexandrov:2021dyl}.  The anti-holomorphic spin fields obey the
same equations with every field and coordinate replaced by its
anti-holomorphic counterpart.  This identical left-right form is another
place where the type IIB chirality choice enters.

\subsection{Conventions for the \texorpdfstring{$bc$}{bc} ghost system}
\label{sec:bc-system}

The reparametrization ghosts $b(z)$ and $c(z)$ are Grassmann odd.  Their
holomorphic conformal weights and ghost numbers are
\begin{center}
\begin{tabular}{ccc}
\toprule
field & $h$ & ghost number \\
\midrule
$b$ & $2$ & $-1$ \\
$c$ & $-1$ & $+1$ \\
\bottomrule
\end{tabular}
\end{center}
\medskip
Their singular operator product is
\begin{equation}
b(z)c(w)\sim\frac{1}{z-w}.
\label{eq:bc-data}
\end{equation}
The $bc$ stress tensor is defined by
\begin{equation}
T_{bc}:=-2b\partial c-(\partial b)c,
\label{eq:Tbc-definition}
\end{equation}
and its central charge is
\begin{equation}
c_{bc}=-26.
\label{eq:cbc}
\end{equation}

On the complex plane, a nonzero $bc$ correlator contains three more $c$
insertions than $b$ insertions.  We define the
constant $K$ by the normalization used in Ref.~\cite{Barman:2026qzi}:
\begin{equation}
\left\langle c(z_1)c(z_2)c(z_3)\right\rangle_{bc}
:=-K(z_1-z_2)(z_2-z_3)(z_1-z_3).
\label{eq:ccc-normalization}
\end{equation}
Here the subscript $bc$ means that the expectation value is taken only in
the $bc$ sector.  Equation~\eqref{eq:ccc-normalization} defines $K$; its
relation to a D-brane or D-instanton tension additionally depends on the
normalization of the string amplitude and is not assumed in these notes.

\subsection{Conventions for the \texorpdfstring{$\beta\gamma$}{beta-gamma} system}
\label{sec:beta-gamma}

The superconformal ghosts $\beta(z)$ and $\gamma(z)$ are Grassmann even and
have conformal weights $3/2$ and $-1/2$, respectively.  We bosonize them in
terms of Grassmann-odd fields $\xi,\eta$ and a Grassmann-even chiral scalar
$\phi$ by
\begin{align}
\beta
&:=\partial\xi\,e^{-\phi},
\label{eq:beta-bosonization}
\\
\gamma
&:=\eta\,e^{\phi}.
\label{eq:gamma-bosonization}
\end{align}
The order of the factors in Eqs.~\eqref{eq:beta-bosonization} and
\eqref{eq:gamma-bosonization} is part of the convention.  The elementary
operator products are
\begin{align}
\xi(z)\eta(w)
&\sim\frac{1}{z-w},
\label{eq:xi-eta-ope}
\\
\partial\phi(z)\partial\phi(w)
&\sim-\frac{1}{(z-w)^2},
\label{eq:dphi-dphi-ope}
\\
e^{q_1\phi}(z)e^{q_2\phi}(w)
&\sim(z-w)^{-q_1q_2}e^{(q_1+q_2)\phi}(w).
\label{eq:exp-phi-ope}
\end{align}
It follows from Eqs.~\eqref{eq:beta-bosonization}--\eqref{eq:exp-phi-ope}
that
\begin{equation}
\beta(z)\gamma(w)\sim-\frac{1}{z-w},
\qquad
\gamma(z)\beta(w)\sim\frac{1}{z-w}.
\label{eq:beta-gamma-ope}
\end{equation}

The bosonized stress tensors are
\begin{align}
T_{\xi\eta}
&:=-\eta\partial\xi,
\label{eq:T-xi-eta}
\\
T_\phi
&:=-\frac12(\partial\phi)^2-\partial^2\phi,
\label{eq:T-phi}
\\
T_{\beta\gamma}
&:=T_{\xi\eta}+T_\phi.
\label{eq:T-beta-gamma}
\end{align}
In the unbosonized variables, the same stress tensor is
\begin{equation}
T_{\beta\gamma}
=-\frac32\beta\partial\gamma
-\frac12(\partial\beta)\gamma.
\label{eq:T-beta-gamma-unbosonized}
\end{equation}
Its central charge is
\begin{equation}
c_{\beta\gamma}=11.
\label{eq:c-beta-gamma}
\end{equation}

The ghost number and picture number used below are additive charges defined
on the elementary bosonized fields by the following table.
\begin{center}
\begin{tabular}{lccccc}
\toprule
field & $\beta$ & $\gamma$ & $\xi$ & $\eta$ & $e^{q\phi}$ \\
\midrule
ghost number & $-1$ & $+1$ & $-1$ & $+1$ & $0$ \\
picture number & $0$ & $0$ & $+1$ & $-1$ & $q$ \\
\bottomrule
\end{tabular}
\end{center}
\medskip
The conformal weight of the exponential is
\begin{equation}
h\!\left[e^{q\phi}\right]
=-\frac12q(q+2).
\label{eq:phi-exponential-weight}
\end{equation}
The mode expansions of the weight-zero field $\xi$ and the weight-one field
$\eta$ define
\begin{equation}
\xi_0:=\oint\frac{\d z}{2\pi\i z}\,\xi(z),
\qquad
\eta_0:=\oint\frac{\d z}{2\pi\i}\,\eta(z).
\label{eq:xi-eta-zero-modes}
\end{equation}
The large Hilbert space is the state space on which $\xi_0$ acts without the
restriction in Eq.~\eqref{eq:small-hilbert-space}.  The small Hilbert space is
the subspace of states $|\Psi\rangle$ satisfying
\begin{equation}
\eta_0|\Psi\rangle=0.
\label{eq:small-hilbert-space}
\end{equation}
Equations~\eqref{eq:beta-bosonization}--\eqref{eq:small-hilbert-space}
follow the ordering, charges, and Hilbert-space convention of
Ref.~\cite{Barman:2026qzi}.

\subsection{Conventions for the \texorpdfstring{$TT$ and $T_FT_F$}{TT and TF TF} operator products}
\label{sec:super-virasoro}

Let $T(z)$ be the stress tensor of a holomorphic matter superconformal field
theory, let $\TF(z)$ be its supercurrent in the normalization defined below,
and let $c_{\mathrm m}$ be its central charge.  We define the normalization
of $T$ and $\TF$ by the singular operator products
\begin{align}
T(z)T(w)
&\sim
\frac{c_{\mathrm m}/2}{(z-w)^4}
+\frac{2T(w)}{(z-w)^2}
+\frac{\partial T(w)}{z-w},
\label{eq:TT-ope}
\\
T(z)\TF(w)
&\sim
\frac{3\TF(w)}{2(z-w)^2}
+\frac{\partial\TF(w)}{z-w},
\label{eq:T-TF-ope}
\\
\TF(z)\TF(w)
&\sim
\frac{c_{\mathrm m}/6}{(z-w)^3}
+\frac{T(w)}{2(z-w)}.
\label{eq:TF-TF-ope}
\end{align}
Thus $\TF$ is one half of a supercurrent $G$ normalized by
$G(z)G(w)\sim(2c_{\mathrm m}/3)(z-w)^{-3}+2T(w)(z-w)^{-1}$.
This explicit algebraic relation is the definition of the relative
normalization; no separate meaning is assigned to the phrase ``normalized
supercurrent.''  Equations~\eqref{eq:TT-ope}--\eqref{eq:TF-TF-ope} match the
convention used in Ref.~\cite{Barman:2026qzi}.

For the ten free bosons and ten free fermions of
Secs.~\ref{sec:free-boson} and \ref{sec:free-fermion},
\begin{equation}
c_{\mathrm m}=c_X+c_\psi=15,
\qquad
T=\Tmat.
\label{eq:critical-matter-central-charge}
\end{equation}
The operator products therefore become
\begin{align}
\Tmat(z)\Tmat(w)
&\sim
\frac{15/2}{(z-w)^4}
+\frac{2\Tmat(w)}{(z-w)^2}
+\frac{\partial\Tmat(w)}{z-w},
\label{eq:critical-TT-ope}
\\
\Tmat(z)\TF(w)
&\sim
\frac{3\TF(w)}{2(z-w)^2}
+\frac{\partial\TF(w)}{z-w},
\label{eq:critical-T-TF-ope}
\\
\TF(z)\TF(w)
&\sim
\frac{5/2}{(z-w)^3}
+\frac{\Tmat(w)}{2(z-w)}.
\label{eq:critical-TF-TF-ope}
\end{align}

For completeness, the coefficients in Eq.~\eqref{eq:critical-TF-TF-ope}
can be reproduced directly.  Substituting
Eq.~\eqref{eq:TF-free-definition}, the double contraction of the two
fermions and the two differentiated bosons gives
\begin{equation}
\left(-\frac{\eta^{MN}}{2(z-w)}\right)
\left(-\frac{\eta_{MN}}{2(z-w)^2}\right)
=\frac{10}{4(z-w)^3}
=\frac{5}{2(z-w)^3}.
\label{eq:TF-double-contraction}
\end{equation}
The single fermion contraction gives
\begin{equation}
-\frac{1}{2(z-w)}\partial X^M\partial X_M(w)
=\frac{T_X(w)}{2(z-w)}.
\label{eq:TF-fermion-single-contraction}
\end{equation}
For the single boson contraction, expand
$\psi^M(z)=\psi^M(w)+(z-w)\partial\psi^M(w)+\cO((z-w)^2)$.
The zeroth-order product $\psi^M\psi_M$ vanishes, while the
first-order term gives
\begin{equation}
-\frac{1}{2(z-w)}(\partial\psi^M)\psi_M(w)
=\frac{T_\psi(w)}{2(z-w)}.
\label{eq:TF-boson-single-contraction}
\end{equation}
Adding Eqs.~\eqref{eq:TF-double-contraction}--\eqref{eq:TF-boson-single-contraction}
gives Eq.~\eqref{eq:critical-TF-TF-ope} with
$\Tmat=T_X+T_\psi$.

The total holomorphic stress tensor of the gauge-fixed critical theory is
\begin{equation}
\Ttot:=\Tmat+T_{bc}+T_{\beta\gamma}.
\label{eq:total-stress-tensor}
\end{equation}
Its central charge vanishes because
\begin{equation}
c_{\mathrm{tot}}
=c_{\mathrm m}+c_{bc}+c_{\beta\gamma}
=15-26+11
=0.
\label{eq:total-central-charge}
\end{equation}

\subsection{Conventions for the BRST and PCO operators}
\label{sec:brst-pco}
\enlargethispage{2\baselineskip}

The holomorphic BRST current is defined by
\begin{equation}
\jB
:=
c\bigl(\Tmat+T_{\xi\eta}+T_\phi\bigr)
+bc\,\partial c
+\gamma\TF
-\frac14\gamma^2b.
\label{eq:brst-current}
\end{equation}
In bosonized variables, the composite field $\gamma^2$ is
\begin{equation}
\gamma^2:=\gamma\gamma
=\eta\,\partial\eta\,e^{2\phi}.
\label{eq:gamma-squared}
\end{equation}
The BRST charge is the contour integral
\begin{equation}
\QB:=\oint\frac{\d z}{2\pi\i}\,\jB(z),
\label{eq:brst-charge}
\end{equation}
where the contour is counterclockwise around the operator on which $\QB$
acts.  The balance $c_{\mathrm m}=15$ in
Eq.~\eqref{eq:critical-matter-central-charge} implies
\begin{equation}
\QB^2=0.
\label{eq:brst-nilpotence}
\end{equation}
The anti-holomorphic charge $\widetilde Q_{\mathrm B}$ is defined by the
tilded copy of Eqs.~\eqref{eq:brst-current}--\eqref{eq:brst-charge}; the
closed-string BRST charge is
\begin{equation}
Q_{\mathrm{closed}}:=\QB+\widetilde Q_{\mathrm B}.
\label{eq:closed-brst-charge}
\end{equation}

For two Grassmann-odd operators $A$ and $B$, write
$\{A,B\}:=AB+BA$.  The holomorphic picture-changing operator (PCO) is
defined in the large Hilbert space by
\begin{align}
\XPCO(z)
&:=2\{\QB,\xi(z)\}
\label{eq:pco-definition}
\\
&=2c\,\partial\xi
+2e^\phi\TF
-\frac12\left[
(\partial\eta)e^{2\phi}b
+\partial\bigl(\eta e^{2\phi}b\bigr)
\right](z).
\label{eq:pco-expanded}
\end{align}
The factor of two in Eq.~\eqref{eq:pco-definition} is part of the convention
adopted from Ref.~\cite{Barman:2026qzi}.  Define
$X_{\mathrm{common}}:=\{\QB,\xi\}$ to be the PCO without this factor.  The
two conventions obey
\begin{equation}
\XPCO=2X_{\mathrm{common}}.
\label{eq:pco-convention-comparison}
\end{equation}
The PCO $\XPCO$ has conformal weight zero, ghost number zero, picture number
$+1$, and is BRST closed:
\begin{equation}
[\QB,\XPCO(z)]=0.
\label{eq:pco-brst-closed}
\end{equation}
The anti-holomorphic PCO is
\begin{equation}
\widetilde\XPCO(\bar z)
:=2\{\widetilde Q_{\mathrm B},\widetilde\xi(\bar z)\}.
\label{eq:antiholomorphic-pco}
\end{equation}

For an oriented worldsheet of genus $g\in\mathbb{Z}_{\geq0}$ with
$b\in\mathbb{Z}_{\geq0}$ boundary components, define $P_{\mathrm{tot}}$ to
be the sum of the picture numbers of all holomorphic and anti-holomorphic
insertions, including PCOs.  A nonzero correlator requires
\begin{equation}
P_{\mathrm{tot}}=2(2g+b-2).
\label{eq:picture-balance}
\end{equation}
For a closed worldsheet, $b=0$, this condition can be imposed separately as
$P_{\mathrm L}=P_{\mathrm R}=2g-2$.  Equation~\eqref{eq:picture-balance}
explains why PCO insertions are required when the picture numbers of the
vertex operators do not already have the required sum.

\begingroup
\small
\bibliographystyle{apsrev4-1long}
\bibliography{references}
\endgroup

\end{document}
